% --- Содержимое файла materials/5_3.tex ---

\subsection{Продолжение оператора по непрерывности}

\begin{idea}[Смысл теоремы]
	Часто бывает сложно задать оператор сразу на всем «сложном» пространстве (например, $L_2$). Гораздо проще задать его на «хорошем» плотном подмножестве (например, на многочленах или непрерывных функциях), а затем «растянуть» на все остальное пространство с помощью предельного перехода.
	
	Эта теорема гарантирует, что такое «растягивание» возможно и корректно, \textbf{только если} пространство, куда мы попадаем ($E_2$), является \textbf{полным} (Банаховым). Иначе пределов последовательностей может просто не существовать.
\end{idea}

\begin{theorem}[О продолжении по непрерывности]\label{extrapolate-to-closure}
	Пусть:
	\begin{itemize}
		\item $E_1$ --- линейное нормированное пространство;
		\item $E_2$ --- \textbf{банахово} пространство (полное);
		\item $D(A) \subset E_1$ --- линейное подпространство, всюду плотное в $E_1$ ($\overline{D(A)} = E_1$);
		\item $A: D(A) \to E_2$ --- линейный ограниченный оператор.
	\end{itemize}
	
	Тогда существует \textbf{единственное} продолжение оператора $A$ до линейного непрерывного оператора $\widetilde{A}$, определенного на всем $E_1$, причем их нормы совпадают:
	\[ \exists! \widetilde{A} \in \mathcal{L}(E_1, E_2): \quad \widetilde{A}\big|_{D(A)} = A \quad \text{и} \quad \|\widetilde{A}\| = \|A\|. \]
\end{theorem}

\begin{proof}
	\textbf{1. Единственность.}
	Пусть существуют два таких непрерывных расширения $\widetilde{A}_1$ и $\widetilde{A}_2$.
	Возьмем произвольный $x \in E_1$. В силу плотности $D(A)$, существует последовательность $\{x_n\} \subset D(A)$, сходящаяся к $x$.
	В силу непрерывности операторов:
	\[ \widetilde{A}_1 x = \lim_{n \to \infty} \widetilde{A}_1 x_n \overset{x_n \in D(A)}{=} \lim_{n \to \infty} A x_n = \lim_{n \to \infty} \widetilde{A}_2 x_n = \widetilde{A}_2 x. \]
	Значит, $\widetilde{A}_1 = \widetilde{A}_2$.
	
	\textbf{2. Существование (Конструкция).}
	Для любого $x \in E_1$ выберем последовательность $\{x_n\} \subset D(A)$, такую что $x_n \to x$.
	Определим значение оператора формулой:
	\[ \widetilde{A}x := \lim_{n \to \infty} A x_n \]
	
	Чтобы это определение было корректным, нужно проверить три факта:
	
	\textbf{а) Предел существует.}
	Так как $x_n \to x$, последовательность $\{x_n\}$ фундаментальна в $E_1$.
	В силу ограниченности (непрерывности) исходного оператора $A$:
	\[ \|Ax_n - Ax_m\| = \|A(x_n - x_m)\| \leqslant \|A\| \cdot \|x_n - x_m\| \to 0. \]
	Значит, последовательность образов $\{Ax_n\}$ фундаментальна в $E_2$.
	Так как $E_2$ --- \textbf{банахово}, эта последовательность имеет предел.
	
	
	
	\textbf{б) Предел не зависит от выбора последовательности.}
	Пусть $x_n \to x$ и $x'_n \to x$. Рассмотрим смешанную последовательность $y = \{x_1, x'_1, x_2, x'_2, \dots\}$.
	Очевидно, $y_k \to x$. Согласно пункту (а), последовательность $\{Ay_k\}$ сходится, а значит, все её подпоследовательности (в частности $\{Ax_n\}$ и $\{Ax'_n\}$) сходятся к одному и тому же пределу.
	
	\textbf{в) Совпадение на $D(A)$.}
	Если $x \in D(A)$, то в качестве последовательности можно взять стационарную $x_n = x$.
	Тогда $\widetilde{A}x = \lim Ax = Ax$.
	
	\textbf{3. Линейность.}
	Следует из линейности $A$ и свойств пределов:
	\[ \widetilde{A}(\alpha x + \beta y) = \lim A(\alpha x_n + \beta y_n) = \alpha \lim Ax_n + \beta \lim Ay_n = \alpha \widetilde{A}x + \beta \widetilde{A}y. \]
	
	\textbf{4. Сохранение нормы.}
	\begin{itemize}
		\item Оценка сверху ($\|\widetilde{A}\| \leqslant \|A\|$):
		\[ \|\widetilde{A}x\| = \lim_{n \to \infty} \|Ax_n\| \leqslant \lim_{n \to \infty} \|A\| \cdot \|x_n\| = \|A\| \cdot \|x\|. \]
		Следовательно, $\|\widetilde{A}\| \leqslant \|A\|$.
		\item Оценка снизу ($\|\widetilde{A}\| \geqslant \|A\|$):
		Так как $\widetilde{A}$ является расширением $A$, то супремум на всем пространстве не может быть меньше супремума на подмножестве:
		\[ \|\widetilde{A}\| = \sup_{x \in E_1, \|x\|=1} \|\widetilde{A}x\| \geqslant \sup_{x \in D(A), \|x\|=1} \|\widetilde{A}x\| = \sup_{x \in D(A), \|x\|=1} \|Ax\| = \|A\|. \]
	\end{itemize}
	Итого: $\|\widetilde{A}\| = \|A\|$.
\end{proof}